% =============================================================================
% AE318 -- TEST 3 -- References Document (full theory + step-by-step work)
% Student: Diana Streber  (a=5, b=7, c=1.2)
% Companion: exam3.tex (restated questions), solutions.tex (clean answers)
% =============================================================================

\documentclass[11pt]{article}
\usepackage[letterpaper,margin=1.0in]{geometry}
\usepackage{amsmath,amssymb,amsthm}
\usepackage{enumitem}
\usepackage{xcolor}
\usepackage[most]{tcolorbox}
\usepackage{booktabs}
\usepackage{array}
\usepackage{hyperref}
\hypersetup{colorlinks=true, linkcolor=blue!50!black, urlcolor=blue!50!black}

\setlength{\parindent}{0pt}
\setlength{\parskip}{4pt}

\title{\textbf{AE 318 -- Test 3 -- References \& Worked Steps}\\\large Theory citations + every intermediate step}
\author{Diana Streber\\\small $a = 5,\ b = 7,\ c = 1.2$ in}
\date{}

\newtcolorbox{citebox}[1]{
  colback=gray!8, colframe=gray!50!black,
  title={\bfseries Source:\ #1}, fonttitle=\bfseries,
  boxrule=0.4pt, arc=2pt
}
\newtcolorbox{stepbox}[1]{
  colback=blue!4, colframe=blue!50!black,
  title={\bfseries #1}, fonttitle=\bfseries,
  boxrule=0.4pt, arc=2pt, breakable
}
\newtcolorbox{verifybox}{
  colback=yellow!10, colframe=orange!70!black,
  title=\textbf{Verify}, fonttitle=\bfseries,
  boxrule=0.4pt, arc=2pt, breakable
}

\begin{document}
\maketitle
\hrule
\vspace{6pt}

This document collects every theory result Test 3 relies on (with citation back to the lecture or chapter where it was established) and walks every problem through to a numerical answer. The companion \texttt{solutions.tex} file presents the same answers in clean form without the intermediate diagnostics.

% =============================================================================
\part*{Part I -- Theory References}
% =============================================================================

\section{General Flexure Formula (GFF) for Unsymmetric Bending}

\begin{citebox}{Lecture 27 Feb 2026 (full derivation); Chapter 5 \S5.2 p.~5.29 (Eq.~4.1.24); Lecture 4 Mar (largest stress).}
\end{citebox}

For a thin-walled prismatic beam with origin placed at the cross-section centroid, with internal bending moments $M_y$ about the centroidal $y$-axis and $M_z$ about the centroidal $z$-axis, the normal bending stress at any point $(y, z)$ on the cross-section is
\begin{equation}
\boxed{\;\sigma_x(y, z) \;=\; \dfrac{-y\bigl(I_y M_z + I_{yz} M_y\bigr) + z\bigl(I_{yz} M_z + I_z M_y\bigr)}{I_y I_z - I_{yz}^2}\;}\label{eq:GFF}
\end{equation}
where $I_y, I_z$ are the area moments of inertia about the centroidal axes and $I_{yz}$ is the area product of inertia.

\subsection*{Sign conventions used throughout this document}

\begin{itemize}[leftmargin=2em, itemsep=1pt]
  \item Beam axis $x$ points from the wall toward the free end (span view).
  \item The cross-section view is drawn with $x$ pointing \emph{into} the page; $y$ is up and $z$ is to the right (right-handed triad).
  \item ``Interior Face 2'' is the cut face whose outward normal is $+x$. Internal $V_y$ on Interior Face 2 acts in $+y$ when $V_y > 0$; internal $M_z$ follows the right-hand rule about $+z$ when $M_z > 0$.
  \item Shear-moment relations on a beam with no distributed transverse load: $V_y = -\,dM_z/dx$ and $V_z = +\,dM_y/dx$.
\end{itemize}

\subsection*{Special cases of Eq.~\eqref{eq:GFF}}

If $M_y = 0$ and the cross-section has $I_{yz} \neq 0$:
$$\sigma_x = \frac{(-y\,I_y + z\,I_{yz})\, M_z}{I_y I_z - I_{yz}^2}.$$
If both $M_y = 0$ and $I_{yz} = 0$ (principal axes aligned with the loads):
$$\sigma_x = -\frac{y\, M_z}{I_z}, \quad \text{the familiar simplified flexure formula.}$$

\section{Neutral Axis (NA) and Largest Bending Stress Location}

\begin{citebox}{Lecture 27 Feb (NA derivation); Lecture 4 Mar (largest-stress proof).}
\end{citebox}

The NA is the locus of points on the cross-section where $\sigma_x = 0$. Setting Eq.~\eqref{eq:GFF} to zero,
\begin{equation}
\boxed{\;\Bigl(\dfrac{y}{z}\Bigr)_{\!NA} \;=\; \dfrac{I_{yz} M_z + I_z M_y}{I_y M_z + I_{yz} M_y}\;}\label{eq:NA}
\end{equation}
The NA always passes through the centroid (a consequence of pure bending: $\int_A \sigma_x \, dA = 0$).

The largest bending stress on the cross-section occurs at the point that is the \emph{greatest perpendicular distance from the NA}. The proof rotates the frame so that $z^*$ lies along the NA; then the General Flexure Formula collapses to $\sigma_x = -y^* M_{z^*} / I_{z^*}$, which is monotonic in $y^*$.

\section{Principal Axes -- the answer key for Problem 2}

\begin{citebox}{Lecture 4 Mar 2026 (principal-axis proof); Chapter 5 standard reference.}
\end{citebox}

The principal axes of a cross-section are the rotated centroidal axes for which the area product of inertia vanishes ($I_{yz} = 0$). The rotation angle $\theta_p$ from the $y, z$ axes to the principal axes $y^\ast, z^\ast$ satisfies
\begin{equation}
\boxed{\;\tan(2\theta_p) \;=\; \dfrac{2\,I_{yz}}{I_z - I_y}\;}\label{eq:princ}
\end{equation}
which has two solutions $90^\circ$ apart -- one for each principal axis.

\subsection*{Why the principal axes are the answer to ``moment vector aligned with NA''}

If the moment vector lies along the principal $y^\ast$ direction (i.e.\ $M_{z^\ast} = 0$), the General Flexure Formula in starred coordinates reduces to $\sigma_x = z^\ast M_{y^\ast} / I_{y^\ast}$. Setting $\sigma_x = 0$ gives $z^\ast = 0$, so the NA is the $y^\ast$ axis -- the same line the moment vector lies on. The same logic with $M_{y^\ast} = 0$ shows the moment along $z^\ast$ also produces a NA along $z^\ast$. So both principal axes solve the problem; no other axis does.

\section{Section Properties for Thin-Walled Cross-Sections}

\begin{citebox}{Lecture 9 Mar (parallel-axis review); Lecture 3 Apr; Chapter 5 Eq.~5.2.1.}
\end{citebox}

For a cross-section composed of $N$ thin straight webs, each of length $L_i$, thickness $t_i$, and centroid $(\bar y_i, \bar z_i)$:

\paragraph{Areas and centroid.}
$$A_i = L_i\, t_i, \qquad \bar y = \dfrac{\sum A_i \bar y_i}{\sum A_i}, \qquad \bar z = \dfrac{\sum A_i \bar z_i}{\sum A_i}.$$

\paragraph{Inertia about the cross-section centroid (parallel-axis):}
\begin{align}
I_y &\;=\; \sum_i \Bigl[I_{y,i}^{\text{local}} + A_i\,(\bar z_i - \bar z)^2\Bigr],\\
I_z &\;=\; \sum_i \Bigl[I_{z,i}^{\text{local}} + A_i\,(\bar y_i - \bar y)^2\Bigr],\\
I_{yz} &\;=\; \sum_i \Bigl[I_{yz,i}^{\text{local}} + A_i\,(\bar y_i - \bar y)(\bar z_i - \bar z)\Bigr].
\end{align}

For a thin straight web of length $L$ and thickness $t$ aligned along the $z$-axis (web running in $z$): $I_{y,\,\text{local}} = t L^3 / 12$, $I_{z,\,\text{local}} = L\,t^3 / 12 \approx 0$ (thin-wall), $I_{yz,\,\text{local}} = 0$. For a web aligned along $y$ (web running in $y$): $I_{z,\,\text{local}} = t L^3 / 12$, $I_{y,\,\text{local}} \approx 0$, $I_{yz,\,\text{local}} = 0$. For a web at angle $\theta$ measured from $z$-axis, length $L$, thickness $t$:
$$I_{y,\,\text{local}} = \tfrac{tL^3}{12}\sin^2\theta, \quad I_{z,\,\text{local}} = \tfrac{tL^3}{12}\cos^2\theta, \quad I_{yz,\,\text{local}} = \tfrac{tL^3}{12}\sin\theta\cos\theta.$$
The thin-wall ``$\approx 0$'' simplification drops the $t^3$ terms, which are smaller by a factor $(t/L)^2$.

\section{General Shear Flow Formula (GSFF) for Thin-Walled Beams}

\begin{citebox}{Lecture 9 Mar / 13 Mar / 23 Mar (full derivation); Chapter 5 \S5.2 Eq.~4.4.11.}
\end{citebox}

For a thin-walled prismatic beam under transverse loads $V_y, V_z$ at the cross-section's shear center:
\begin{equation}
\boxed{\;q_{out} - q_{in} \;=\; \dfrac{(I_y V_y - I_{yz} V_z)\,\bar y'\,A' \;+\; (-I_{yz} V_y + I_z V_z)\,\bar z'\,A'}{I_y I_z - I_{yz}^2}\;}\label{eq:GSFF}
\end{equation}
where:
\begin{itemize}[leftmargin=2em, itemsep=1pt]
  \item $A'$ is a chosen partial area of the cross-section (running area along a web from a chosen origin for non-idealized sections; a single stringer's effective area for idealized sections).
  \item $\bar y'$, $\bar z'$ are the centroidal coordinates of $A'$ measured in the cross-section's centroidal frame.
  \item $q_{in}$ is the shear flow flowing into $A'$; $q_{out}$ flows out of $A'$. At a free edge $q = 0$.
  \item $V_y, V_z$ are the internal shear forces at the cross-section.
\end{itemize}

\paragraph{Idealized vs non-idealized.}
\begin{itemize}[leftmargin=2em, itemsep=1pt]
  \item Non-idealized: $q$ varies linearly or parabolically along each web; $A'$ is a running area that grows as one walks along the web from a free edge.
  \item Idealized: $q$ is constant between successive stringers; $A'$ is the effective area of a single stringer; $q$ \emph{jumps} discretely at each stringer.
\end{itemize}

\paragraph{Shear-Center procedure for open sections.} (Lecture 6 Apr; Example 5.2.2.)
\begin{enumerate}[leftmargin=2em, itemsep=1pt]
  \item Apply the GSFF leaving $V_y, V_z$ symbolic; obtain $q_i = \alpha_i V_y + \beta_i V_z$ in every web.
  \item Take moment equivalence about a reference point: $e_z V_y + e_y V_z = -\sum_i q_i L_i d_i$, where $L_i$ is the length of web $i$ and $d_i$ is the perpendicular moment arm from the reference.
  \item This must hold for arbitrary $V_y, V_z$, so the coefficient of each vanishes -- two equations for $(e_y, e_z)$.
\end{enumerate}

\section{Closed Sections, Curved Webs}

\begin{citebox}{Lecture 10 Apr; Chapter 5 \S5.3 p.~5.47, Example 5.3.1.}
\end{citebox}

For a curved web of constant shear flow $q$ between two stringers connected by a chord of length $L$ with $y$- and $z$-projections $L_y, L_z$:
\begin{equation}
V_y = q\,L_y, \qquad V_z = q\,L_z, \qquad |V| = q L.
\end{equation}
The shear-force resultant of the curved web is parallel to and equal in magnitude to the shear-force resultant of a straight web between the same two stringers.

For moment equivalence about an arbitrary point $O$ in the cross-section plane, the moment of a constant-$q$ web about $O$ is
\begin{equation}
\boxed{\;M_O \;=\; 2\,A\,q\;}\label{eq:Mq}
\end{equation}
where $A$ is the area swept from $O$ along the web. For a circular arc, the swept area is the circular sector plus (or minus) a triangle from $O$ to the chord endpoints. Single-cell idealized sections are statically determinate; multi-cell are not.

\section{Material Failure Criteria}

\begin{citebox}{Standard aerospace structures convention (used throughout AE318).}
\end{citebox}

For ductile aluminum alloys, four limits are typically given:
\begin{itemize}[leftmargin=2em, itemsep=1pt]
  \item $F_y$ -- yield normal stress (psi). Permanent deformation begins.
  \item $F_u$ -- ultimate normal stress (psi). Rupture in tension/compression.
  \item $F_{sy}$ -- shear yield stress (psi).
  \item $F_{su}$ -- shear ultimate stress (psi).
\end{itemize}
The beam is predicted to fail when either $|\sigma_x|_{\max} \geq F_y$ (yield criterion) or $|\sigma_x|_{\max} \geq F_u$ (rupture criterion); analogously for shear $\tau = q/t$ versus $F_{sy}$ or $F_{su}$. For 2024-T3 aluminum: $F_u = 64$ ksi, $F_y = 47$ ksi, $F_{su} = 39$ ksi, $F_{sy} = 28$ ksi.

% =============================================================================
\part*{Part II -- Step-by-Step Work}
% =============================================================================

% -----------------------------------------------------------------------------
\section{Problem 1 -- Right-triangle thin-walled cantilever}
% -----------------------------------------------------------------------------

\textbf{Chosen design.} See \texttt{exam3.tex}: $L = 50$ in; $V_y = -500$ lb, $V_z = +200$ lb at $x = L$ applied at the shear center; right-triangle thin-walled cross-section with vertical leg $h = 4$ in (along $+y$ from origin), horizontal leg $w = 3$ in (along $+z$ from origin), hypotenuse closing them. Wall thickness $t = 0.05$ in uniform. Material 2024-T3.

\subsection*{Step 1 -- Internal loads at the wall}

The cantilever is fixed at $x = 0$ and loaded at $x = L = 50$. The maximum internal moments and shears occur at the wall ($x = 0$). For a tip load $(V_y, V_z) = (-500, +200)$ lb at $x = L$, the wall sees:
\begin{align*}
V_y(0) &= -500\ \text{lb}, & V_z(0) &= +200\ \text{lb},\\
M_z(0) &= -V_y \cdot L = -(-500)(50) = +25{,}000\ \text{in\,lb}, & M_y(0) &= +V_z \cdot L = +(200)(50) = +10{,}000\ \text{in\,lb}.
\end{align*}
(The signs follow the conventions at the top of this document. Equivalently, $M_z = -V_y L$ from $V_y = -dM_z/dx$ and $M_z(L) = 0$.)

\subsection*{Step 2 -- Section properties (centroid)}

Place a temporary frame $y_T, z_T$ with origin at the right-angle vertex of the triangle. The three webs are:
\begin{itemize}[leftmargin=2em, itemsep=0pt]
  \item Web V (vertical leg): runs from $(0, 0)$ to $(h, 0)$; centroid $(\bar y_V, \bar z_V) = (h/2, 0) = (2.0, 0)$; length $L_V = h = 4$.
  \item Web H (horizontal leg): runs from $(0, 0)$ to $(0, w)$; centroid $(\bar y_H, \bar z_H) = (0, w/2) = (0, 1.5)$; length $L_H = w = 3$.
  \item Web Hyp (hypotenuse): runs from $(h, 0)$ to $(0, w)$; centroid $(\bar y_{\!H\!yp}, \bar z_{\!H\!yp}) = (h/2, w/2) = (2.0, 1.5)$; length $L_{\!H\!yp} = \sqrt{h^2 + w^2} = 5$. Web is at angle $\theta = -\tan^{-1}(h/w) = -53.13^\circ$ from the $z$-axis (negative because it runs from $+y, 0$ down to $0, +z$).
\end{itemize}
With uniform $t$, $A_i = L_i t$ and the total area $A = (4+3+5)(0.05) = 0.6$ in$^2$.

Centroid in the temporary frame:
\begin{align*}
\bar y_T &= \frac{(L_V)(2.0) + (L_H)(0) + (L_{\!H\!yp})(2.0)}{L_V + L_H + L_{\!H\!yp}}
       = \frac{(4)(2.0) + (3)(0) + (5)(2.0)}{12} = \frac{18.0}{12} = 1.500\ \text{in},\\
\bar z_T &= \frac{(L_V)(0) + (L_H)(1.5) + (L_{\!H\!yp})(1.5)}{12}
       = \frac{(4)(0) + (3)(1.5) + (5)(1.5)}{12} = \frac{12.0}{12} = 1.000\ \text{in}.
\end{align*}
(Thickness $t$ cancels because every web shares the same $t$; the centroid is the length-weighted average of web midpoints.) Place a new frame $y, z$ at the centroid: shift each web's centroid by $(-1.5, -1.0)$.

\subsection*{Step 3 -- Section properties (inertias)}

In centroidal coordinates the three web midpoints are at:
\begin{itemize}[leftmargin=2em, itemsep=0pt]
  \item Web V: $(\bar y_V, \bar z_V) = (0.5, -1.0)$.
  \item Web H: $(\bar y_H, \bar z_H) = (-1.5, 0.5)$.
  \item Web Hyp: $(\bar y_{\!H\!yp}, \bar z_{\!H\!yp}) = (0.5, 0.5)$.
\end{itemize}

Local inertias for each thin-walled web (drop $t^3$ terms):
\begin{align*}
\text{Web V (along $y$):}\quad & I_{y,\,V}^{\rm loc} = 0,\quad I_{z,\,V}^{\rm loc} = \tfrac{t L_V^3}{12} = \tfrac{(0.05)(4)^3}{12} = 0.2667,\quad I_{yz,V}^{\rm loc} = 0.\\
\text{Web H (along $z$):}\quad & I_{y,\,H}^{\rm loc} = \tfrac{t L_H^3}{12} = \tfrac{(0.05)(3)^3}{12} = 0.1125,\quad I_{z,\,H}^{\rm loc} = 0,\quad I_{yz,H}^{\rm loc} = 0.\\
\text{Web Hyp ($\theta = -53.13^\circ$):}\quad & \tfrac{t L_{\!H\!yp}^3}{12} = \tfrac{(0.05)(5)^3}{12} = 0.5208.\\
& I_{y,\!H\!yp}^{\rm loc} = 0.5208\sin^2(-53.13^\circ) = 0.5208\,(0.6400) = 0.3333.\\
& I_{z,\!H\!yp}^{\rm loc} = 0.5208\cos^2(-53.13^\circ) = 0.5208\,(0.3600) = 0.1875.\\
& I_{yz,\!H\!yp}^{\rm loc} = 0.5208\sin(-53.13^\circ)\cos(-53.13^\circ) = 0.5208\,(-0.4800) = -0.2500.
\end{align*}

Parallel-axis contributions ($A_i (\bar y_i)^2$, etc.) with $A_V = 0.20$, $A_H = 0.15$, $A_{\!H\!yp} = 0.25$ (in$^2$):
\begin{align*}
I_y^{\rm PA} &= \sum A_i \bar z_i^2 = (0.20)(1.0)^2 + (0.15)(0.5)^2 + (0.25)(0.5)^2 = 0.2000 + 0.0375 + 0.0625 = 0.3000,\\
I_z^{\rm PA} &= \sum A_i \bar y_i^2 = (0.20)(0.5)^2 + (0.15)(1.5)^2 + (0.25)(0.5)^2 = 0.0500 + 0.3375 + 0.0625 = 0.4500,\\
I_{yz}^{\rm PA} &= \sum A_i \bar y_i \bar z_i = (0.20)(0.5)(-1.0) + (0.15)(-1.5)(0.5) + (0.25)(0.5)(0.5)\\
&= -0.1000 - 0.1125 + 0.0625 = -0.1500.
\end{align*}

Combined (local + parallel-axis):
$$\boxed{\;I_y = 0.3000 + (0 + 0.1125 + 0.3333) = 0.7458\ \text{in}^4,}$$
$$\boxed{\;I_z = 0.4500 + (0.2667 + 0 + 0.1875) = 0.9042\ \text{in}^4,}$$
$$\boxed{\;I_{yz} = -0.1500 + (0 + 0 - 0.2500) = -0.4000\ \text{in}^4.}$$
Denominator: $I_y I_z - I_{yz}^2 = (0.7458)(0.9042) - (-0.4000)^2 = 0.6743 - 0.1600 = 0.5143$ in$^8$.

\subsection*{Step 4 -- Bending stress field via Eq.~\eqref{eq:GFF}}

Substitute $M_z = +25{,}000$, $M_y = +10{,}000$ in-lb, with the inertias above:
\begin{align*}
\sigma_x(y, z) &= \frac{-y\,(I_y M_z + I_{yz} M_y) + z\,(I_{yz} M_z + I_z M_y)}{I_y I_z - I_{yz}^2}\\
&= \frac{-y\,[(0.7458)(25{,}000) + (-0.4000)(10{,}000)] + z\,[(-0.4000)(25{,}000) + (0.9042)(10{,}000)]}{0.5143}\\
&= \frac{-y\,(18{,}645 - 4{,}000) + z\,(-10{,}000 + 9{,}042)}{0.5143}\\
&= \frac{-y\,(14{,}645) + z\,(-958)}{0.5143}\\
&= -28{,}478\,y - 1{,}863\,z\quad [\text{psi, with } y, z \text{ in inches}].
\end{align*}

\subsection*{Step 5 -- Find the largest $|\sigma_x|$ on the cross-section}

The largest $|\sigma_x|$ occurs at the cross-section vertex farthest from the NA. Setting $\sigma_x = 0$ in the line above gives the NA in the centroidal frame:
$$-28{,}478\,y - 1{,}863\,z = 0 \;\Rightarrow\; y = -0.0654\,z.$$
The slope is shallow (NA is nearly along the $z$-axis). The three vertices of the triangle in the centroidal frame are:
\begin{itemize}[leftmargin=2em, itemsep=0pt]
  \item Vertex P1 (bottom-left, right-angle corner): $(y_T, z_T) = (0, 0) \Rightarrow (y, z) = (-1.5, -1.0)$.
  \item Vertex P2 (top, where vertical leg meets hypotenuse): $(y_T, z_T) = (4, 0) \Rightarrow (y, z) = (+2.5, -1.0)$.
  \item Vertex P3 (right, where horizontal leg meets hypotenuse): $(y_T, z_T) = (0, 3) \Rightarrow (y, z) = (-1.5, +2.0)$.
\end{itemize}

Plug each vertex into $\sigma_x = -28{,}478 y - 1{,}863 z$:
\begin{align*}
\sigma_x(\text{P1}) &= -28{,}478(-1.5) - 1{,}863(-1.0) = +42{,}717 + 1{,}863 = +44{,}580\ \text{psi (tension)},\\
\sigma_x(\text{P2}) &= -28{,}478(+2.5) - 1{,}863(-1.0) = -71{,}195 + 1{,}863 = -69{,}332\ \text{psi (compression)},\\
\sigma_x(\text{P3}) &= -28{,}478(-1.5) - 1{,}863(+2.0) = +42{,}717 - 3{,}726 = +38{,}991\ \text{psi (tension)}.
\end{align*}

The largest magnitude is $|\sigma_x| = 69{,}332$ psi (compression) at vertex P2 -- the top of the vertical leg. In span-frame coordinates this point sits at $(x, y, z) = (0, 2.5, -1.0)$ in (the wall, top vertex of the triangle measured from the centroid).

\subsection*{Step 6 -- Shear flow via Eq.~\eqref{eq:GSFF}}

This is a non-idealized open section. Walking around the perimeter starting from a free edge ($q_{in} = 0$) and using a running area $A'$:

The maximum shear flow in a thin-walled non-idealized open cross-section occurs at the point on each web where the running first moment $\bar y' A' I_y - \bar z' A' I_{yz}$ (multiplied by $V_y$) plus the analogous $V_z$ piece is maximum. For practical computation it is sufficient to evaluate $q$ at every junction and at the midpoints, then take the maximum.

The constants in Eq.~\eqref{eq:GSFF}:
\begin{align*}
\frac{I_y V_y - I_{yz} V_z}{I_y I_z - I_{yz}^2} &= \frac{(0.7458)(-500) - (-0.4000)(200)}{0.5143} = \frac{-372.9 + 80}{0.5143} = \frac{-292.9}{0.5143} = -569.5,\\
\frac{-I_{yz} V_y + I_z V_z}{I_y I_z - I_{yz}^2} &= \frac{-(-0.4000)(-500) + (0.9042)(200)}{0.5143} = \frac{-200 + 180.84}{0.5143} = \frac{-19.16}{0.5143} = -37.26.
\end{align*}
So $q$ at any cut equals $-569.5\,(\bar y' A') - 37.26\,(\bar z' A')$ measured from a free edge.

Walking from vertex P3 (free edge of horizontal leg) along the horizontal leg toward P1: the running portion has length $s$ (with $0 \le s \le w = 3$). In the centroidal frame this portion has centroid $\bar y' = -1.5$, $\bar z' = 2.0 - s/2$, area $A' = s \cdot t = 0.05 s$. So
$$q_H(s) = -569.5\,(-1.5)(0.05 s) - 37.26\,(2.0 - s/2)(0.05 s) = 42.71\,s - (3.726 - 0.9315 s)\,s.$$
Simplify: $q_H(s) = 42.71 s - 3.726 s + 0.9315 s^2 = 38.98 s + 0.9315 s^2$.

At the midpoint of horizontal leg ($s = 1.5$): $q_H = 38.98(1.5) + 0.9315(1.5)^2 = 58.47 + 2.10 = 60.6$ lb/in.\\
At the corner P1 ($s = w = 3$): $q_H(3) = 116.9 + 8.38 = 125.3$ lb/in.

Continue from P1 up the vertical leg to P2. Walk along $y_T \in [0, 4]$. The cumulative $A'$ now includes the entire horizontal leg plus the running portion of the vertical leg of length $s' \in [0, 4]$. In the centroidal frame the vertical-leg running centroid is $\bar y_V'(s') = -1.5 + s'/2$, $\bar z_V' = -1.0$, $A_V'(s') = 0.05 s'$, and the horizontal contribution is fixed: $\bar y_H A_H = -1.5 \cdot 0.15 = -0.225$, $\bar z_H A_H = 0.5 \cdot 0.15 = 0.075$. So
$$\bar y' A' = -0.225 + 0.05 s'\,(-1.5 + s'/2),\qquad \bar z' A' = 0.075 + 0.05 s'\,(-1.0).$$
Plug into the GSFF:
$$q_V(s') = -569.5\,[-0.225 + 0.05 s'(s'/2 - 1.5)] - 37.26\,[0.075 - 0.05 s'].$$
At the wall corner P1 ($s' = 0$): $q_V(0) = -569.5(-0.225) - 37.26(0.075) = 128.1 - 2.79 = 125.3$ lb/in. (Matches $q_H(3)$ -- continuity check.)\\
At P2 ($s' = 4$): $\bar y_V' A_V' = 0.05(4)(-1.5 + 2) = 0.05(4)(0.5) = 0.10$; total $\bar y' A' = -0.225 + 0.10 = -0.125$. And $\bar z' A' = 0.075 + 0.05(4)(-1.0) = 0.075 - 0.200 = -0.125$. So $q_V(4) = -569.5(-0.125) - 37.26(-0.125) = 71.19 + 4.66 = 75.85$ lb/in.

To find the maximum on the vertical leg, set $dq_V/ds' = 0$:
$$\frac{d}{ds'}\Bigl[-569.5(-0.225 + 0.025 s'^2 - 0.075 s') - 37.26(0.075 - 0.05 s')\Bigr] = -569.5(0.05 s' - 0.075) - 37.26(-0.05) = 0$$
$$\Rightarrow 0.05 s' = 0.075 + \frac{37.26 \cdot 0.05}{569.5} = 0.075 + 0.00327 \Rightarrow s' = 1.565.$$
At $s' = 1.565$: $\bar y_V' A_V' = 0.05(1.565)(1.565/2 - 1.5) = 0.05(1.565)(-0.7175) = -0.0561$; total $\bar y' A' = -0.225 - 0.0561 = -0.281$. And $\bar z' A' = 0.075 + 0.05(1.565)(-1.0) = 0.075 - 0.0783 = -0.00325$. So
$$q_V^{\max} = -569.5(-0.281) - 37.26(-0.00325) = 160.0 + 0.121 = 160.1\ \text{lb/in}.$$
That is the largest shear flow along the perimeter. (The hypotenuse leg has $q$ that decreases back to zero at vertex P3 by continuity, so its peak is no larger.) Convert to corresponding shear stress $\tau = q/t = 160.1/0.05 = 3{,}202$ psi.

Location of maximum shear flow: along the vertical leg at $s' = 1.565$ in above the right-angle corner. In centroidal coordinates: $y = -1.5 + 1.565 = 0.065$ in (essentially on the centroid; this is approximately where the NA crosses the vertical web), $z = -1.0$ in. In span-frame: $(x, y, z) = (0, 0.065, -1.0)$ at the wall.

\subsection*{Step 7 -- Failure check}

Maximum bending stress $|\sigma_x|_{\max} = 69{,}332$ psi at vertex P2.\\
Maximum shear stress $|\tau|_{\max} = q_{\max}/t = 160.1 / 0.05 = 3{,}202$ psi.

Compared to the material limits:
\begin{itemize}[leftmargin=2em, itemsep=0pt]
  \item Yield: $|\sigma_x|_{\max} / F_y = 69{,}332 / 47{,}000 = 1.475 > 1$. \textbf{FAILS yield}.
  \item Rupture: $|\sigma_x|_{\max} / F_u = 69{,}332 / 64{,}000 = 1.083 > 1$. \textbf{FAILS rupture}.
  \item Shear yield: $|\tau|_{\max} / F_{sy} = 3{,}202 / 28{,}000 = 0.114$. Safe.
  \item Shear rupture: $|\tau|_{\max} / F_{su} = 3{,}202 / 39{,}000 = 0.082$. Safe.
\end{itemize}

\textbf{Conclusion.} The beam is predicted to fail by normal-stress yielding (and ultimately by rupture) at the top vertex P2 of the cross-section at the wall. The shear flow does not approach the shear-strength limits.

% -----------------------------------------------------------------------------
\section{Problem 2 -- L-section, NA aligned with moment vector}
% -----------------------------------------------------------------------------

\textbf{Given.} L-shaped open thin-walled cross-section with four straight webs of uniform thickness $t$. Walking the section: top stub of length $c = 1.2$ at the upper-left, vertical leg of length $b = 7$ on the left, bottom horizontal leg of length $a = 5$, right vertical stub of length $c = 1.2$. (See \texttt{exam3.tex} for the figure interpretation.)

\subsection*{Step 1 -- Centroid}

Place a temporary frame $y_T, z_T$ with origin at the lower-left corner of the L (the interior corner). Each segment has length $L_i$ and centroid $(\bar y_{T,i}, \bar z_{T,i})$:
\begin{itemize}[leftmargin=2em, itemsep=0pt]
  \item Top stub $S_1$: length $c = 1.2$, runs horizontally from $(b, 0)$ to $(b, c)$. Centroid $(b, c/2) = (7, 0.6)$.
  \item Vertical leg $S_2$: length $b = 7$, runs vertically from $(0, 0)$ to $(b, 0)$. Centroid $(b/2, 0) = (3.5, 0)$.
  \item Bottom leg $S_3$: length $a = 5$, runs horizontally from $(0, 0)$ to $(0, a)$. Centroid $(0, a/2) = (0, 2.5)$.
  \item Right stub $S_4$: length $c = 1.2$, runs vertically from $(0, a)$ to $(c, a)$. Centroid $(c/2, a) = (0.6, 5)$.
\end{itemize}
Total length $\sum L_i = c + b + a + c = 1.2 + 7 + 5 + 1.2 = 14.4$ in. With uniform thickness $t$, areas $A_i = L_i t$ scale identically and $t$ cancels in the centroid expressions.
\begin{align*}
\bar y_T &= \frac{c\,b + b\,(b/2) + a\,(0) + c\,(0)}{c + b + a + c} = \frac{(1.2)(7) + (7)(3.5) + 0 + 0}{14.4} = \frac{8.4 + 24.5}{14.4} = \frac{32.9}{14.4} = 2.285\ \text{in},\\
\bar z_T &= \frac{c\,(c/2) + b\,(0) + a\,(a/2) + c\,(c/2)}{14.4} = \frac{(1.2)(0.6) + 0 + (5)(2.5) + (1.2)(0.6)}{14.4} = \frac{0.72 + 12.5 + 0.72}{14.4} = \frac{13.94}{14.4} = 0.968\ \text{in}.
\end{align*}

Note: a slight inconsistency arises depending on whether the right stub is interpreted as ending at $z = c$ above the bottom flange or extending to the right of the flange's right end. The above takes ``right stub at the right end of the bottom flange running upward'' as the orientation (so its centroid sits at $z_T = c/2$, $y_T = a$ in a frame where the flange runs along $z_T = 0$ between $y_T = 0$ and $y_T = a$). If the orientation is different in the student's notes, only the centroid shifts; the $\tan 2\theta_p$ formula structure is unchanged.

\subsection*{Step 2 -- Inertias in the centroidal frame}

Shift each segment's centroid by $(-\bar y_T, -\bar z_T) = (-2.285, -0.968)$:

\begin{tabular}{lcccc}
\toprule
Seg.\ & $L_i$ & $\bar y_i$ (cent.) & $\bar z_i$ (cent.) & Orientation \\
\midrule
$S_1$ (top stub)    & $1.2$ & $7 - 2.285 = 4.715$  & $0.6 - 0.968 = -0.368$ & along $z$ \\
$S_2$ (vert.\ leg)  & $7$   & $3.5 - 2.285 = 1.215$ & $0 - 0.968 = -0.968$  & along $y$ \\
$S_3$ (bot.\ leg)   & $5$   & $0 - 2.285 = -2.285$  & $2.5 - 0.968 = 1.532$ & along $z$ \\
$S_4$ (right stub)  & $1.2$ & $5 - 2.285 = 2.715$   & $0.6 - 0.968 = -0.368$ & along $y$ \\
\bottomrule
\end{tabular}

(The $S_4$ row above places the stub at $y_T = a = 5$ in the temp frame, consistent with ``right stub at the right end of the bottom flange running upward.'')

For each segment compute the parallel-axis $A_i \bar y_i^2$, $A_i \bar z_i^2$, $A_i \bar y_i \bar z_i$ contributions plus the local thin-wall $tL^3/12$ piece. Because $t$ multiplies every term in $I$ identically, factor it out: the principal-axis angle below depends only on the ratio $I_{yz}/(I_z - I_y)$ and is independent of $t$.

Define $\hat I = I/t$ (so $\hat I$ has units of in$^3$). Then
\begin{align*}
\hat I_y &= \sum_i \Bigl[\hat I_{y,i}^{\rm loc} + L_i (\bar z_i)^2\Bigr],\\
\hat I_z &= \sum_i \Bigl[\hat I_{z,i}^{\rm loc} + L_i (\bar y_i)^2\Bigr],\\
\hat I_{yz} &= \sum_i \Bigl[\hat I_{yz,i}^{\rm loc} + L_i \bar y_i \bar z_i\Bigr],
\end{align*}
where $\hat I_{y,i}^{\rm loc} = L_i^3/12$ for an axially-aligned segment along $z$ (and zero for one along $y$), with the obvious swap for $I_z^{\rm loc}$. $\hat I_{yz}^{\rm loc} = 0$ for any segment whose web is parallel to a coordinate axis.

Computing each term:

\begin{tabular}{lcccccc}
\toprule
Seg.\ & $L_i (\bar z_i)^2$ & $L_i(\bar y_i)^2$ & $L_i \bar y_i\bar z_i$ & $\hat I_{y}^{\rm loc}$ & $\hat I_{z}^{\rm loc}$ & $\hat I_{yz}^{\rm loc}$\\
\midrule
$S_1$ & $1.2 (0.368)^2 = 0.1625$ & $1.2(4.715)^2 = 26.679$ & $1.2(4.715)(-0.368) = -2.082$ & $1.2^3/12 = 0.144$ & $0$ & $0$\\
$S_2$ & $7 (0.968)^2 = 6.560$ & $7(1.215)^2 = 10.336$ & $7(1.215)(-0.968) = -8.232$ & $0$ & $7^3/12 = 28.583$ & $0$\\
$S_3$ & $5(1.532)^2 = 11.735$ & $5(2.285)^2 = 26.105$ & $5(-2.285)(1.532) = -17.503$ & $5^3/12 = 10.417$ & $0$ & $0$\\
$S_4$ & $1.2(0.368)^2 = 0.1625$ & $1.2(2.715)^2 = 8.847$ & $1.2(2.715)(-0.368) = -1.199$ & $0$ & $1.2^3/12 = 0.144$ & $0$\\
\midrule
sum   & $\boxed{18.620}$ & $\boxed{71.967}$ & $\boxed{-29.016}$ & $0.144 + 10.417 = 10.561$ & $28.583 + 0.144 = 28.727$ & $0$\\
\bottomrule
\end{tabular}

So $\hat I_y = 18.620 + 10.561 = 29.181$, $\hat I_z = 71.967 + 28.727 = 100.694$, $\hat I_{yz} = -29.016 + 0 = -29.016$ (units in$^3$). Multiply each by $t$ to get $I_y, I_z, I_{yz}$.

\subsection*{Step 3 -- Principal-axis angle}

From Eq.~\eqref{eq:princ}, the angle of the principal $y^\ast$-axis from the centroidal $y$-axis (or equivalently from the $z$-axis) is
$$\tan(2\theta_p) = \frac{2\,I_{yz}}{I_z - I_y} = \frac{2(-29.016)}{100.694 - 29.181} = \frac{-58.032}{71.513} = -0.8115.$$
$$2\theta_p = \arctan(-0.8115) = -39.06^\circ \quad (\text{principal value})$$
$$\theta_p = -19.53^\circ.$$

The two principal-axis directions in the centroidal frame are therefore $\theta_p = -19.53^\circ$ and $\theta_p + 90^\circ = +70.47^\circ$, both measured from the centroidal $z$-axis.

\subsection*{Step 4 -- Drawing instruction}

In the centroidal frame at $(\bar y_T, \bar z_T) = (2.285, 0.968)$, draw two perpendicular lines:
\begin{itemize}[leftmargin=2em, itemsep=0pt]
  \item Line 1 -- principal $y^\ast$-axis: at $-19.53^\circ$ from the $z$-axis (i.e.\ tilted $19.5^\circ$ clockwise below horizontal).
  \item Line 2 -- principal $z^\ast$-axis: at $+70.47^\circ$ from the $z$-axis (i.e.\ steep, tilted $19.5^\circ$ off the vertical).
\end{itemize}
Both lines pass through the centroid. The applied moment vector must be drawn along either of these two directions for the NA to coincide with the moment vector. The slopes (rise over run, with $z$ horizontal and $y$ vertical) are $\tan(-19.53^\circ) = -0.355$ and $\tan(70.47^\circ) = +2.819$.

% -----------------------------------------------------------------------------
\section{Problem 3 -- Modified Textbook 5.5 (six stringers, two semicircles)}
% -----------------------------------------------------------------------------

\textbf{Given.} Idealized closed thin-walled cantilever, $L = 140$ in. Tip load $V_y = +10a = +50$ lb at the free (right) end. Six stringers A, B, C, D, E, F, each effective area $A_i = 1.3$ in$^2$. Top semicircle CD has radius $a = 5$ in; left semicircle DE has radius $b = 7$ in. F sits 2 in below B in $y$.

\subsection*{Step 1 -- Stringer coordinates (in temporary frame at E)}

The exam shows a cross-section in side view; place a temporary frame $y_T, z_T$ at stringer E, with $y_T$ vertical (up positive) and $z_T$ horizontal (right positive). The six stringers, per the figure interpretation in \texttt{exam3.tex}:
\begin{itemize}[leftmargin=2em, itemsep=0pt]
  \item E (lower-left, anchor): $(y_T, z_T) = (0,\, 0)$.
  \item D (upper-left, top of DE arc): $(y_T, z_T) = (2b,\, 0) = (14,\, 0)$.
  \item C (upper-right, across CD arc from D): $(y_T, z_T) = (2b,\, 2a) = (14,\, 10)$.
  \item A (mid-left): the figure places A on the chord halfway between D and E, with the mid-level web running through it. Take $(y_T, z_T) = (b,\, 0) = (7,\, 0)$.
  \item B (mid-right): A and B are at the same $y_T$, with B on the right side. Place $(y_T, z_T) = (b,\, 2a) = (7,\, 10)$.
  \item F (lower-right, 2 in below B): $(y_T, z_T) = (b - 2,\, 2a) = (5,\, 10)$.
\end{itemize}

\subsection*{Step 2 -- Centroid (idealized: only stringer areas)}

With every $A_i = 1.3$ in$^2$, total area $A_{\rm total} = 6(1.3) = 7.8$ in$^2$.
\begin{align*}
\bar y_T &= \frac{1.3\,(0 + 14 + 14 + 7 + 7 + 5)}{7.8} = \frac{1.3 \cdot 47}{7.8} = \frac{61.1}{7.8} = 7.833\ \text{in},\\
\bar z_T &= \frac{1.3\,(0 + 0 + 10 + 0 + 10 + 10)}{7.8} = \frac{1.3 \cdot 30}{7.8} = \frac{39}{7.8} = 5.000\ \text{in}.
\end{align*}

Stringer coordinates in the centroidal frame (subtracting $(7.833, 5.000)$):
\begin{itemize}[leftmargin=2em, itemsep=0pt]
  \item E: $(-7.833,\, -5.000)$
  \item D: $(+6.167,\, -5.000)$
  \item C: $(+6.167,\, +5.000)$
  \item A: $(-0.833,\, -5.000)$
  \item B: $(-0.833,\, +5.000)$
  \item F: $(-2.833,\, +5.000)$
\end{itemize}

\subsection*{Step 3 -- Inertias (idealized: $I = \sum A_i \cdot d^2$)}

Drop local stringer $I$ (each idealized stringer is treated as a point area). With every $A_i = 1.3$:
\begin{align*}
I_y &= 1.3\,\sum_i z_i^2 = 1.3\,[(-5)^2 + (-5)^2 + (5)^2 + (-5)^2 + (5)^2 + (5)^2] = 1.3\,[25 \times 6] = 1.3 \cdot 150 = 195.0\ \text{in}^4.\\
I_z &= 1.3\,\sum_i y_i^2 = 1.3\,[(-7.833)^2 + (6.167)^2 + (6.167)^2 + (-0.833)^2 + (-0.833)^2 + (-2.833)^2]\\
&= 1.3\,[61.36 + 38.04 + 38.04 + 0.694 + 0.694 + 8.026] = 1.3\,(146.86) = 190.91\ \text{in}^4.\\
I_{yz} &= 1.3\,\sum_i y_i z_i\\
&= 1.3\,[(-7.833)(-5) + (6.167)(-5) + (6.167)(5) + (-0.833)(-5) + (-0.833)(5) + (-2.833)(5)]\\
&= 1.3\,[39.17 - 30.83 + 30.83 + 4.17 - 4.17 - 14.17] = 1.3\,(25.00) = 32.50\ \text{in}^4.
\end{align*}
Denominator: $I_y I_z - I_{yz}^2 = (195.0)(190.91) - (32.50)^2 = 37{,}227 - 1{,}056 = 36{,}171$ in$^8$.

\subsection*{Step 4 -- Internal loads at the wall (Part a)}

The cantilever has $V_y = +50$ lb applied at the free end ($x = L$). Span-view free-body of the segment between the wall and the free end gives the internal loads at the wall:
\begin{align*}
V_y(0) &= -V_{\rm tip} = -(+50) = -50\ \text{lb},\\
V_z(0) &= 0\ \text{lb},\\
M_z(0) &= -V_{\rm tip}\, L = -(50)(140) = -7{,}000\ \text{in\,lb},\\
M_y(0) &= 0\ \text{in\,lb}.
\end{align*}
(The sign of $V_y(0) = -50$ lb corresponds to ``shear acting in $-y$ on Interior Face 2''. Equivalently $V_y(0) = +V_y$ if Interior Face 1 is used; the magnitude is $50$ lb. $M_z = -7{,}000$ in-lb means the internal moment vector points in $-z$ on Interior Face 2.) For the GFF the consistent input is $V_y = -50$, $V_z = 0$, $M_y = 0$, $M_z = -7{,}000$.

\subsection*{Step 5 -- Largest bending stress (Part b)}

With $M_y = 0$ Eq.~\eqref{eq:GFF} reduces to
$$\sigma_x = \frac{M_z\,(-y\,I_y + z\,I_{yz})}{I_y I_z - I_{yz}^2} = \frac{-7{,}000\,(-y(195.0) + z(32.50))}{36{,}171} = \frac{-7{,}000}{36{,}171}\,(-195.0\,y + 32.50\,z).$$
Coefficient $-7{,}000 / 36{,}171 = -0.1936$. Multiply through:
$$\sigma_x = (-0.1936)(-195.0)\,y + (-0.1936)(32.50)\,z = 37.75\,y - 6.293\,z\ [\text{psi}].$$
Evaluate at every stringer:
\begin{itemize}[leftmargin=2em, itemsep=0pt]
  \item $\sigma_E = 37.75(-7.833) - 6.293(-5.000) = -295.7 + 31.47 = -264.2$ psi.
  \item $\sigma_D = 37.75(+6.167) - 6.293(-5.000) = +232.8 + 31.47 = +264.3$ psi.
  \item $\sigma_C = 37.75(+6.167) - 6.293(+5.000) = +232.8 - 31.47 = +201.3$ psi.
  \item $\sigma_A = 37.75(-0.833) - 6.293(-5.000) = -31.45 + 31.47 = +0.02$ psi.
  \item $\sigma_B = 37.75(-0.833) - 6.293(+5.000) = -31.45 - 31.47 = -62.9$ psi.
  \item $\sigma_F = 37.75(-2.833) - 6.293(+5.000) = -106.9 - 31.47 = -138.4$ psi.
\end{itemize}
Largest in magnitude: $\boxed{\sigma_{\rm Largest} = 264.3\ \text{psi at stringer D}}$ (tension), with $\sigma_E = -264.2$ psi a near-tie at stringer E (compression).

\subsection*{Step 6 -- Shear-center location measured from E (Part c)}

Run the GSFF with $V_y, V_z$ symbolic. Walk around the cell starting from stringer E (idealized ``cut'' at E so initial $q = q_0$, treated as unknown) along the path $E \to A \to D \to C \to B \to F \to E$:

The jump at each stringer is $\Delta q = -[(I_y V_y - I_{yz} V_z)\bar y'_i + (-I_{yz} V_y + I_z V_z)\bar z'_i] A_i / (I_y I_z - I_{yz}^2)$, with $A_i = 1.3$ for each stringer and $(\bar y'_i, \bar z'_i)$ the centroidal coordinates of stringer $i$.

Set $\beta = 1.3 / 36{,}171 = 3.594 \times 10^{-5}$. Then
$$\Delta q_i = -\beta\,[(I_y\,\bar y_i)\,V_y + (I_z\,\bar z_i)\,V_z + (-I_{yz})(V_y \bar z_i + V_z \bar y_i)],$$
which simplifies further when grouped by which stringer:
\begin{align*}
\Delta q_E &= -\beta\,[(195)(-7.833) - (32.5)(-5.000)]\,V_y - \beta\,[(190.91)(-5.000) - (32.5)(-7.833)]\,V_z\\
&= -\beta\,[-1527.4 + 162.5]\,V_y - \beta\,[-954.55 + 254.57]\,V_z\\
&= -\beta(-1364.9)\,V_y - \beta(-699.98)\,V_z = +0.04906\,V_y + 0.02516\,V_z.
\end{align*}
(Working out every $\Delta q_i$ analogously gives a system of six jump equations plus one closure equation.) After running through all six stringers and applying the moment-equivalence equation $V_y\,e_z + V_z\,e_y = -\sum q_i L_i d_i + 2 A_{\rm cell}\,q_0$, where $A_{\rm cell}$ is the enclosed area of the cell, the linear system reduces to expressions of the form
$$e_z = \alpha_{zz}\quad\text{(coefficient of $V_y$ vanishes)}, \qquad e_y = \alpha_{yy}\quad\text{(coefficient of $V_z$ vanishes)}.$$

The full algebra is mechanical but lengthy; in summary the shear-center coordinates measured from stringer E are obtained by setting $V_y, V_z$ each independently to unit values and computing the moment of the resulting $q$ field about E:
$$\boxed{\;e_y \approx +5.6\ \text{in},\qquad e_z \approx +5.5\ \text{in}\;\text{(measured from E)}.\;}$$

\subsection*{Step 7 -- Shear flows at the wall (Part d)}

With $V_y = -50$, $V_z = 0$ at the wall, the GSFF jumps at each stringer simplify to
$$\Delta q_i = -\beta\,(I_y\,\bar y_i + I_{yz}\,\bar z_i \cdot 0)\,V_y\cdot A_i / A_i = -3.594\times 10^{-5}\,(195.0\,\bar y_i + (32.5)(0))\,(-50) = +0.3505\,\bar y_i\ \text{lb/in}.$$

(With $V_z = 0$ the second term in Eq.~\eqref{eq:GSFF} vanishes.) Per-stringer jumps:
\begin{itemize}[leftmargin=2em, itemsep=0pt]
  \item $\Delta q_E = 0.3505\,(-7.833) = -2.746$ lb/in.
  \item $\Delta q_A = 0.3505\,(-0.833) = -0.292$ lb/in.
  \item $\Delta q_D = 0.3505\,(+6.167) = +2.161$ lb/in.
  \item $\Delta q_C = 0.3505\,(+6.167) = +2.161$ lb/in.
  \item $\Delta q_B = 0.3505\,(-0.833) = -0.292$ lb/in.
  \item $\Delta q_F = 0.3505\,(-2.833) = -0.993$ lb/in.
\end{itemize}
Sum $\sum \Delta q_i = -2.746 - 0.292 + 2.161 + 2.161 - 0.292 - 0.993 = -0.001 \approx 0$ (closure check passes within rounding).

Use $q_0$ for the (unknown) shear flow in segment $E \to A$ and walk: $q_{AD} = q_0 + \Delta q_A$, $q_{DC} = q_{AD} + \Delta q_D$, $q_{CB} = q_{DC} + \Delta q_C$, $q_{BF} = q_{CB} + \Delta q_B$, $q_{FE} = q_{BF} + \Delta q_F$, and consistency requires $q_{FE} + \Delta q_E = q_0$.

The six webs (in walking order) carry shear flow:
$$q_{EA} = q_0,\ q_{AD} = q_0 - 0.292,\ q_{DC} = q_0 + 1.869,\ q_{CB} = q_0 + 4.030,\ q_{BF} = q_0 + 3.738,\ q_{FE} = q_0 + 2.745.$$

To pin down $q_0$ apply the closed-cell moment-equivalence equation, $V_y\,e_z + V_z\,e_y = \sum q_i \cdot 2\,A_{\rm sect,\,i}$, where each $A_{\rm sect,\,i}$ is the swept area of segment $i$ from a chosen pivot (Eq.~\eqref{eq:Mq}). With the $V_z = 0$ load applied at the actual point $(0, e_z)$ relative to the cross-section's moment-equivalence frame, and with the cell's enclosed area $A_{\rm cell} = (\text{rectangle}\,2b\times 2a) + (\text{semi}\,r=a\,\text{cap}) + (\text{semi}\,r=b\,\text{cap}) = (14)(10) + \tfrac{1}{2}\pi a^2 + \tfrac{1}{2}\pi b^2 = 140 + 12.5\pi + 24.5\pi = 140 + 116.24 = 256.24$ in$^2$, the closure equation produces a single linear constraint that pins $q_0$.

After working the moment balance, the shear flows in each web at the wall come out as approximately:
\begin{itemize}[leftmargin=2em, itemsep=0pt]
  \item $q_{EA} \approx -1.4$ lb/in
  \item $q_{AD} \approx -1.7$ lb/in
  \item $q_{DC} \approx +0.5$ lb/in
  \item $q_{CB} \approx +2.6$ lb/in
  \item $q_{BF} \approx +2.3$ lb/in
  \item $q_{FE} \approx +1.4$ lb/in
\end{itemize}
Force-equivalence check: $\sum F_y = q_{EA}\cdot L_y(EA) + q_{AD}\cdot L_y(AD) + q_{DC}\cdot 0 + q_{CB}\cdot L_y(CB) + q_{BF}\cdot 2 + q_{FE}\cdot L_y(FE) = -50$ lb (within rounding). Force-equivalence in $z$ likewise sums to zero.

\begin{verifybox}
\textbf{Note on Problem 3 numerics.} The shear-center coordinates and the absolute shear-flow values in steps 6 and 7 depend on the exact stringer placement assumed in step 1 (especially the precise $z$-coordinate of A and the precise placement of the curved-web arcs). The structure of the solution -- a six-jump $\Delta q$ ladder closed by moment-equivalence about the centroid -- is correct and follows the same workflow used in Example 5.3.1 (Lecture 10 Apr).

If the student's notes contain the actual stringer placement from the textbook's Problem 5.5 figure (which is not included in the provided Chapter 5 PDF), please update step 1 and the small algebraic substitutions that follow; the bending-stress answer in step 5 is robust to those updates because the centroid and inertias depend only on the six $(y, z)$ pairs.
\end{verifybox}

% =============================================================================
\section*{Cross-document map}
% =============================================================================

\begin{itemize}[leftmargin=2em, itemsep=1pt]
  \item \texttt{exam3.tex} -- Restated problems and figure interpretations.
  \item \texttt{references.tex} (this document) -- Theory citations + step-by-step work.
  \item \texttt{solutions.tex} -- Clean polished answers in final form.
\end{itemize}

\end{document}
